\section{Experimental Setup}

The proposed methodology was evaluated on two publicly available transformer language models using emotion
datasets containing either 9 or 20 emotion categories. All experiments follow the activation 
extraction, steering, and evaluation pipeline described in Section~\ref{sec:methodology}. The implementation
was developed in Python via Google Colab notebooks, using the PyTorch, HuggingFace Transformers, and 
scikit-learn libraries.

\subsection{Models}
Two language models were evaluated:

\begin{itemize}
    \item \texttt{openai-community/gpt-2-medium}, a publicly available 355-million-parameter autoregressive Transformer model released by OpenAI, serving as a reproducible baseline representative of early large language models \cite{radford_language_2019}.
    \item \texttt{google/gemma-4-e2b}, a recently released open-weight language model from Google designed for instruction-following and multilingual reasoning \cite{team_gemma_2026}.
\end{itemize}

These models provide contrasting cases in architecture, scale, tokenizer, training corpus, and multilingual capability. GPT-2 Medium serves as an earlier-generation baseline, whereas Gemma 4 E2B provides a modern instruction-tuned case. Recurring patterns across the two models are treated as qualitative convergence under the shared pipeline, not as sufficient evidence of a general property of all Transformer representations.

\subsection{Emotion Datasets}

The proposed methodology was evaluated using two emotion datasets derived from a common emotion-story corpus. The 
first corresponds to the publicly available nine-emotion list adopted by \cite{hsieh_replicating_2026}, providing a reproducible
baseline for comparison. Building upon this benchmark, we additionally construct an extended dataset containing twenty
emotion categories to evaluate the proposed methodology under a broader and finer-grained affective space.

Unlike \cite{sofroniew_emotion_2026}, which investigated 171 emotional concepts, this work intentionally focuses on smaller
emotion datasets that are practical to analyze and compare while remaining computationally reproducible inside free-tier
computing hosts. The twenty-emotion dataset extends the original nine emotions with semantically related affective states, 
allowing us to inspect how the qualitative geometry changes with a more granular label set. To maintain a consistent 
comparison, the nine-emotion experiments were conducted using the corresponding subset of the twenty-emotion emotional stories,
so that the two settings differ in label-set size rather than in the generation procedure used for their shared categories.

Stories were used as generated, without manual rewriting or post-hoc quality filtering. This was an intentional first-pass design decision that keeps the corpus-construction procedure simple and reproducible, but it also bounds interpretation: the resulting plots and interventions describe the behavior of the present generated corpus and extraction pipeline, rather than a corpus-independent property of either model.

\subsection{Extraction Configuration}

All experiments were conducted using a common extraction configuration across both evaluated language models. Emotion Vectors
were extracted from the residual stream representations described in the previous section, followed by normalization, Logit Lens vocabulary
analysis, principal component visualization, and activation-intervention diagnostics. Identical extraction settings were maintained for the 
9-emotion and 20-emotion story datasets, reducing configuration changes as an explanation for differences between the two label sets. Table~\ref{tab:extraction-config} summarizes the configuration employed 
throughout the experiments. The extraction layer was selected as the layer closest to two-thirds of each model's depth, following the mid-late-layer convention used by \cite{sofroniew_emotion_2026}; we did not conduct a layer sweep or select the layer using the reported outcomes.

\begin{table}[t]
\centering
\small
\renewcommand{\arraystretch}{1.2}

\caption{Extraction Configuration Summary}
\label{tab:extraction-config}

\begin{tabularx}{0.75\linewidth}{
    >{\arraybackslash}p{4.3cm}
    @{\hspace{2cm}}
    >{\raggedright\arraybackslash}X
}
\toprule
\textbf{Parameter} & \textbf{Value} \\
\midrule

Stories per emotion & 100 \\
Neutral texts & 50 \\
Hidden layer & 16* and 23** \\
Hidden-state aggregation & Mean pooling \\
Emotion vector normalization & L2 normalization \\
Logit Lens standardization & $z$-score \\
Top-$k$ tokens & 5 \\
Bottom-$k$ tokens & 5 \\
PCA components & 2 \\

\bottomrule
\end{tabularx}

\vspace{0.4em}

\begin{minipage}{0.75\linewidth}
\footnotesize
\textbf{Note.}
* \texttt{openai-community/gpt-2-medium} uses hidden layer 16 ($\sim\nicefrac{2}{3}$ of layer depth). ** \texttt{google/gemma-4-e2b} uses hidden layer 23 ($\sim\nicefrac{2}{3}$ of layer depth).
\end{minipage}

\end{table}

\subsection{Steering Configuration}

The steering experiments were adapted from the evaluation setting used by Anthropic \cite{sofroniew_emotion_2026}. Two neutral prompt templates (Prompt00 and Prompt01) were
employed to elicit unconstrained text continuations before and after applying emotion vector steering. By minimizing the emotional content of 
the initial prompt, the resulting generations provide a common prompt context for descriptive comparisons of the injected directions' effects on
the generated text. Throughout all steering experiments, a fixed steering coefficient of $\alpha=+0.5$ was used because the source study presents steering results at that strength. In this implementation, the intervention is scaled by the current token's residual norm; therefore, this value is a source-informed operating point, not a claim of optimality or a directly comparable absolute strength across models. Text was generated using stochastic sampling with a temperature of 0.7 and a maximum generation length of 100 new tokens.
Both prompt templates were evaluated for the 9-emotion and 20-emotion datasets across \texttt{openai-community/gpt-2-medium} and 
\texttt{google/gemma-4-e2b}. The generation hyperparameters were held constant, although individual sampled continuations remain stochastic. Table~\ref{tab:steering-config} summarizes the configuration employed throughout 
the experiments.


\begin{table}[t]
\centering
\small
\renewcommand{\arraystretch}{1.2}

\caption{Steering Configuration Summary}
\label{tab:steering-config}

\begin{tabularx}{0.75\linewidth}{
    >{\arraybackslash}p{4.3cm}
    @{\hspace{2cm}}
    >{\raggedright\arraybackslash}X
}
\toprule
\textbf{Parameter} & \textbf{Value} \\
\midrule

Steering coefficient $\alpha$ & +0.5 \\
Prompt templates & Prompt00 and Prompt01 \\
Maximum generated tokens & 100 \\
Temperature & 0.7 \\
Sampling & True \\

\bottomrule
\end{tabularx}

\vspace{0.4em}

\begin{minipage}{0.75\linewidth}
\footnotesize
\textbf{Note.} Both Prompt00 and Prompt01 are available on Appendix ~\ref{apdx:primers}.
\end{minipage}

\end{table}

\subsection{Evaluation Pipeline}

The evaluation followed a fixed pipeline for both \texttt{openai-community/gpt-2-medium} and \texttt{google/gemma-4-e2b}. For each
model, emotion vectors were first extracted independently from the 9-emotion and 20-emotion datasets using the extraction 
configuration described previously. The resulting emotion libraries were then evaluated through four complementary analyses. First,
Logit Lens extraction identified the vocabulary tokens most strongly aligned with each extracted direction. Second,
each direction was injected at the selected transformer block while generating continuations from Prompt00 and Prompt01. Changes in
next-token probabilities for the vector-derived target-token sets were visualized through delta-log-probability heatmaps, and representative
sampled continuations were inspected qualitatively. Third, PCA visualized the leading variance directions without prespecifying valence or
arousal axes. Finally, cosine-similarity heatmaps summarized pairwise directional alignment. Together these analyses provide complementary
descriptions of the extracted vectors; none serves as an independent semantic ground truth.

\subsection{Implementation Details}

All experiments were implemented in Python using Google Colab notebooks. Model inference was performed with the HuggingFace 
\emph{Transformers} library and PyTorch, while activation extraction, residual stream steering, and Logit Lens analyses were
implemented using custom Python scripts. Principal Component Analysis (PCA) was performed with scikit-learn, and quantitative
analyses were conducted using NumPy and Pandas. Visualizations, including PCA projections, Cosine Similarity Heatmaps, and Delta 
Log Probability heatmaps, were generated using Matplotlib and Plotly. To reduce redundant computation, extracted emotion vectors
were cached and reused throughout the evaluation pipeline. Finally, all experiments were executed using the free-tier Google Colab
environment equipped with NVIDIA T4 GPUs.

\subsection{Computational Resources}

Unlike the large-scale infrastructure available to Anthropic, the experiments presented in this work were intentionally
designed to operate within free-tier commodity compute environments. Table ~\ref{tab:experimental-scope} below summarizes the principal 
differences in computational scope.

\begin{table}[!htbp]
\centering
\small
\caption{Experimental Scope Comparison of Anthropic versus Current Manuscript}
\label{tab:experimental-scope}

    \begin{tabularx}{\linewidth}{>{\bfseries}p{2.8cm}XX}
    \toprule
    \textbf{Aspect} & \textbf{Anthropic} & \textbf{Current Manuscript} \\
    \midrule

    \textbf{Models} &
    Claude Sonnet 4.5 &
    GPT 2 Medium and Gemma 4 E2B \\[0.4em]

    \textbf{Emotions} &
    171 &
    9 and 20 \\[0.4em]

    \textbf{Emotion stories available} &
    205,200 &
    2,000 \\[0.4em]

    \textbf{Researchers} &
    16 &
    1 \\[0.4em]

    \textbf{Hardware} &
    Internal compute cluster &
    2 Google Colab Notebooks with Free-Tier T4 NVIDIA GPUs \\

    \bottomrule
    \end{tabularx}

    \vspace{0.3em}
    \begin{minipage}{\linewidth}
        \footnotesize
        \textbf{Note.} Inspired by a similar table found on \cite{hsieh_replicating_2026}.
    \end{minipage}
\end{table}
